\section{Limitations}
\label{sec:limitations}

\subsection{Internal validity}
\label{sec:limits-internal}

Each real window is one realised path. The intervals reported on the real windows (Wilson intervals of the hit rates, bootstrap intervals of the score change) describe the sampling variability of the days within the path, and not the variability of the crisis that produced it, so a window that is a single event cannot support a claim about the frequency with which the wrapper repairs coverage in crises in general; the generator experiments supply the replication, for a data-generating process chosen by the author. The windows are short: the pandemic window holds 104 days and 2 to 13 hits per cell in the unwrapped runs, and the Kupiec and conditional coverage p-values are asymptotic, which the engine documents as crude at counts of this size. A cell that passes or fails by a small margin (a p-value of 0.049 fails and 0.050 passes) is sensitive to one hit more or less; \cref{tab:main-results} shows several such cells, among them the wrapped HS-250 at the 1\% level in the crisis window and at the 2.5\% level in the pandemic window, whose conditional coverage p-values are both 0.055 and whose Kupiec p-values of 0.021 and 0.022 decide the fail. The count of passed cells is therefore a coarse summary, and the hit rates and score changes with their intervals are the better evidence. The hypothesis was written with a rule that counts windows passed; a different threshold of the p-value would change the counts.

The step and the band were fixed before the stress windows were run, and the other recalibrations were given constants fixed before their runs; but the choice of which recalibrations to compare, of the tracker's step grid and of the ablation grid was made after the first results on three models were known, and the sensitivity of the tracker to its step shows that the ranking of methods depends on that choice. The ten-model extension used the same constants as the three-model test. No multiplicity correction is applied to the \tenRuns{} cells.

\subsection{External validity}
\label{sec:limits-external}

The portfolios are public series (the US equity market, a ten-year Treasury position, the euro against the dollar) and three mixed portfolios with constant weights; a bank's portfolio has options, credit spreads and positions that change daily, and the VaR of a bank is computed from a full revaluation or a sensitivity model, not from the return series of a market index. The wrapper needs the model to answer at another level; for a bank's production model that is a commitment on its own, and for the five ladder models fitted for one level it requires a refit that the experiments do not test. The three models without an ES cannot be wrapped at all. The stress windows are the ones that Contraste chose, all of the 2000s and 2020s United States and euro markets; they contain neither a currency crisis of an emerging market nor a sovereign event. The tail levels are 1\% and 2.5\%; the supervisory ES level and the 99\% VaR of the older standard are the first of these, but the behaviour at 0.1\% is not tested, and at that level the band's lower edge of 0.001 leaves no room.

\subsection{Construct validity}
\label{sec:limits-construct}

The pass rule (Kupiec and conditional coverage both at least 0.05) is the product's, not the supervisor's: the Basel traffic light counts exceedances over 250 days and has its own zones \citep{bcbsmar32}: at the 1\% level it is green up to four exceptions and red from ten \citep{bcbs1996}, a reading by count alone that ignores the clustering the independence test sees. The experiments tabulate hits, not zones, so a statement that the wrapped models would have been in a lower zone is not made; and a window shorter than 250 days (the pandemic and the 2023 windows) has no zone at all. A bank that reads its record through the zones would therefore see the effect of the wrapper only through the count of hits that the main table gives. The score tolerance of 5\% is a decision. The FZ0 score is a consistent score for the pair, but it is one score among a family, and a different member could rank the recalibrations differently. The independence test of the engine has a first-order alternative and no power against dependence at longer lags, and the duration-based alternative of \citet{christoffersenpelletier2004} was not run. The ES is tested by the statistic of \cref{eq:z2} only; the more recent ES back-tests \citep{duescanciano2017,kratz2018,bayer2022,wang2022} were not run on the wrapped pair, and a wrapped ES may fail them for reasons that $Z_2$ does not see.

\subsection{Scope of the claims}
\label{sec:limits-scope}

The manuscript claims that, on ten ladder models and three stress windows, a wrapper with a pre-set step passes more cells than the unwrapped models without a score cost over 5\% and without an independence failure, that it holds on known-truth generators of persistent and transient shifts, and that its repair is bounded by the band and by the scale of the error. It does not claim the update as a novelty (it is the update of \citet{gibbs2021}); it does not claim to be better than the other recalibrations compared, and the results show that it is not better by cells passed; it does not claim to repair the ES or the independence; and it does not claim that its guarantee, which is vacuous at these window lengths, explains the empirical repair. The reproduction by an independent group on other portfolios is the test the claim has not had. The French equity series is distributed for derived use, so the equity returns cannot be redistributed with the code, and a reader needs the series from the library to rerun the equity windows.
